\element{proportions}{
\begin{question}{prop-partie-1}
Calculer $5\%$ de $80$.
\begin{choiceshoriz}
\correctchoice{$4$}
\wrongchoice{$85$}
\wrongchoice{$16$}
\wrongchoice{$40$}
\end{choiceshoriz}
\end{question}
}

\element{proportions}{
\begin{question}{prop-partie-2}
Calculer $10\%$ de $100$.
\begin{choiceshoriz}
\correctchoice{$10$}
\wrongchoice{$110$}
\wrongchoice{$100$}
\wrongchoice{$1$}
\end{choiceshoriz}
\end{question}
}

\element{proportions}{
\begin{question}{prop-partie-3}
Calculer $15\%$ de $120$.
\begin{choiceshoriz}
\correctchoice{$18$}
\wrongchoice{$135$}
\wrongchoice{$8$}
\wrongchoice{$180$}
\end{choiceshoriz}
\end{question}
}

\element{proportions}{
\begin{question}{prop-partie-4}
Calculer $20\%$ de $140$.
\begin{choiceshoriz}
\correctchoice{$28$}
\wrongchoice{$160$}
\wrongchoice{$7$}
\wrongchoice{$280$}
\end{choiceshoriz}
\end{question}
}

\element{proportions}{
\begin{question}{prop-partie-5}
Calculer $25\%$ de $160$.
\begin{choiceshoriz}
\correctchoice{$40$}
\wrongchoice{$185$}
\wrongchoice{$6{,}4$}
\wrongchoice{$400$}
\end{choiceshoriz}
\end{question}
}

\element{proportions}{
\begin{question}{prop-partie-6}
Calculer $30\%$ de $180$.
\begin{choiceshoriz}
\correctchoice{$54$}
\wrongchoice{$210$}
\wrongchoice{$6$}
\wrongchoice{$540$}
\end{choiceshoriz}
\end{question}
}

\element{proportions}{
\begin{question}{prop-partie-7}
Calculer $35\%$ de $200$.
\begin{choiceshoriz}
\correctchoice{$70$}
\wrongchoice{$235$}
\wrongchoice{$5{,}71$}
\wrongchoice{$700$}
\end{choiceshoriz}
\end{question}
}

\element{proportions}{
\begin{question}{prop-partie-8}
Calculer $40\%$ de $220$.
\begin{choiceshoriz}
\correctchoice{$88$}
\wrongchoice{$260$}
\wrongchoice{$5{,}5$}
\wrongchoice{$880$}
\end{choiceshoriz}
\end{question}
}

\element{proportions}{
\begin{question}{prop-partie-9}
Calculer $45\%$ de $240$.
\begin{choiceshoriz}
\correctchoice{$108$}
\wrongchoice{$285$}
\wrongchoice{$5{,}33$}
\wrongchoice{$1080$}
\end{choiceshoriz}
\end{question}
}

\element{proportions}{
\begin{question}{prop-partie-10}
Calculer $50\%$ de $260$.
\begin{choiceshoriz}
\correctchoice{$130$}
\wrongchoice{$310$}
\wrongchoice{$5{,}2$}
\wrongchoice{$1300$}
\end{choiceshoriz}
\end{question}
}

\element{proportions}{
\begin{question}{prop-total-1}
$12$ représente $10\%$ d'une quantité. Quelle est cette quantité ?
\begin{choiceshoriz}
\correctchoice{$120$}
\wrongchoice{$1{,}2$}
\wrongchoice{$22$}
\wrongchoice{$12$}
\end{choiceshoriz}
\end{question}
}

\element{proportions}{
\begin{question}{prop-total-2}
$15$ représente $15\%$ d'une quantité. Quelle est cette quantité ?
\begin{choiceshoriz}
\correctchoice{$100$}
\wrongchoice{$2{,}25$}
\wrongchoice{$30$}
\wrongchoice{$10$}
\end{choiceshoriz}
\end{question}
}

\element{proportions}{
\begin{question}{prop-total-3}
$18$ représente $20\%$ d'une quantité. Quelle est cette quantité ?
\begin{choiceshoriz}
\correctchoice{$90$}
\wrongchoice{$3{,}6$}
\wrongchoice{$38$}
\wrongchoice{$9$}
\end{choiceshoriz}
\end{question}
}

\element{proportions}{
\begin{question}{prop-total-4}
$21$ représente $25\%$ d'une quantité. Quelle est cette quantité ?
\begin{choiceshoriz}
\correctchoice{$84$}
\wrongchoice{$5{,}25$}
\wrongchoice{$46$}
\wrongchoice{$8{,}4$}
\end{choiceshoriz}
\end{question}
}

\element{proportions}{
\begin{question}{prop-total-5}
$24$ représente $30\%$ d'une quantité. Quelle est cette quantité ?
\begin{choiceshoriz}
\correctchoice{$80$}
\wrongchoice{$7{,}2$}
\wrongchoice{$54$}
\wrongchoice{$8$}
\end{choiceshoriz}
\end{question}
}

\element{proportions}{
\begin{question}{prop-total-6}
$27$ représente $35\%$ d'une quantité. Quelle est cette quantité ?
\begin{choiceshoriz}
\correctchoice{$77{,}14$}
\wrongchoice{$9{,}45$}
\wrongchoice{$62$}
\wrongchoice{$7{,}71$}
\end{choiceshoriz}
\end{question}
}

\element{proportions}{
\begin{question}{prop-total-7}
$30$ représente $10\%$ d'une quantité. Quelle est cette quantité ?
\begin{choiceshoriz}
\correctchoice{$300$}
\wrongchoice{$3$}
\wrongchoice{$40$}
\wrongchoice{$30$}
\end{choiceshoriz}
\end{question}
}

\element{proportions}{
\begin{question}{prop-total-8}
$33$ représente $15\%$ d'une quantité. Quelle est cette quantité ?
\begin{choiceshoriz}
\correctchoice{$220$}
\wrongchoice{$4{,}95$}
\wrongchoice{$48$}
\wrongchoice{$22$}
\end{choiceshoriz}
\end{question}
}

\element{proportions}{
\begin{question}{prop-total-9}
$36$ représente $20\%$ d'une quantité. Quelle est cette quantité ?
\begin{choiceshoriz}
\correctchoice{$180$}
\wrongchoice{$7{,}2$}
\wrongchoice{$56$}
\wrongchoice{$18$}
\end{choiceshoriz}
\end{question}
}

\element{proportions}{
\begin{question}{prop-total-10}
$39$ représente $25\%$ d'une quantité. Quelle est cette quantité ?
\begin{choiceshoriz}
\correctchoice{$156$}
\wrongchoice{$9{,}75$}
\wrongchoice{$64$}
\wrongchoice{$15{,}6$}
\end{choiceshoriz}
\end{question}
}

\element{proportions}{
\begin{question}{prop-quatrieme-1}
Dans un tableau de proportionnalité, $2$ correspond à $5$. À quoi correspond $6$ ?
\begin{choiceshoriz}
\correctchoice{$15$}
\wrongchoice{$2{,}4$}
\wrongchoice{$9$}
\wrongchoice{$13$}
\end{choiceshoriz}
\end{question}
}

\element{proportions}{
\begin{question}{prop-quatrieme-2}
Dans un tableau de proportionnalité, $3$ correspond à $6$. À quoi correspond $9$ ?
\begin{choiceshoriz}
\correctchoice{$18$}
\wrongchoice{$4{,}5$}
\wrongchoice{$12$}
\wrongchoice{$1$}
\end{choiceshoriz}
\end{question}
}

\element{proportions}{
\begin{question}{prop-quatrieme-3}
Dans un tableau de proportionnalité, $4$ correspond à $7$. À quoi correspond $12$ ?
\begin{choiceshoriz}
\correctchoice{$21$}
\wrongchoice{$6{,}86$}
\wrongchoice{$15$}
\wrongchoice{$23$}
\end{choiceshoriz}
\end{question}
}

\element{proportions}{
\begin{question}{prop-quatrieme-4}
Dans un tableau de proportionnalité, $5$ correspond à $8$. À quoi correspond $15$ ?
\begin{choiceshoriz}
\correctchoice{$24$}
\wrongchoice{$9{,}38$}
\wrongchoice{$18$}
\wrongchoice{$28$}
\end{choiceshoriz}
\end{question}
}

\element{proportions}{
\begin{question}{prop-quatrieme-5}
Dans un tableau de proportionnalité, $6$ correspond à $9$. À quoi correspond $18$ ?
\begin{choiceshoriz}
\correctchoice{$27$}
\wrongchoice{$12$}
\wrongchoice{$21$}
\wrongchoice{$33$}
\end{choiceshoriz}
\end{question}
}

\element{proportions}{
\begin{question}{prop-quatrieme-6}
Dans un tableau de proportionnalité, $7$ correspond à $10$. À quoi correspond $21$ ?
\begin{choiceshoriz}
\correctchoice{$30$}
\wrongchoice{$14{,}7$}
\wrongchoice{$24$}
\wrongchoice{$38$}
\end{choiceshoriz}
\end{question}
}

\element{proportions}{
\begin{question}{prop-quatrieme-7}
Dans un tableau de proportionnalité, $8$ correspond à $11$. À quoi correspond $24$ ?
\begin{choiceshoriz}
\correctchoice{$33$}
\wrongchoice{$17{,}45$}
\wrongchoice{$27$}
\wrongchoice{$43$}
\end{choiceshoriz}
\end{question}
}

\element{proportions}{
\begin{question}{prop-quatrieme-8}
Dans un tableau de proportionnalité, $9$ correspond à $12$. À quoi correspond $27$ ?
\begin{choiceshoriz}
\correctchoice{$36$}
\wrongchoice{$20{,}25$}
\wrongchoice{$30$}
\wrongchoice{$48$}
\end{choiceshoriz}
\end{question}
}

\element{proportions}{
\begin{question}{prop-quatrieme-9}
Dans un tableau de proportionnalité, $10$ correspond à $13$. À quoi correspond $30$ ?
\begin{choiceshoriz}
\correctchoice{$39$}
\wrongchoice{$23{,}08$}
\wrongchoice{$33$}
\wrongchoice{$53$}
\end{choiceshoriz}
\end{question}
}

\element{proportions}{
\begin{question}{prop-quatrieme-10}
Dans un tableau de proportionnalité, $11$ correspond à $14$. À quoi correspond $33$ ?
\begin{choiceshoriz}
\correctchoice{$42$}
\wrongchoice{$25{,}93$}
\wrongchoice{$36$}
\wrongchoice{$58$}
\end{choiceshoriz}
\end{question}
}

\element{proportions}{
\begin{question}{prop-echelle-1}
Sur un plan à l'échelle $1:1000$, $2$ cm représentent quelle distance réelle ?
\begin{choiceshoriz}
\correctchoice{$2000$ cm}
\wrongchoice{$1002$ cm}
\wrongchoice{$20$ cm}
\wrongchoice{$500$ cm}
\end{choiceshoriz}
\end{question}
}

\element{proportions}{
\begin{question}{prop-echelle-2}
Sur un plan à l'échelle $1:2000$, $3$ cm représentent quelle distance réelle ?
\begin{choiceshoriz}
\correctchoice{$6000$ cm}
\wrongchoice{$2003$ cm}
\wrongchoice{$60$ cm}
\wrongchoice{$666{,}67$ cm}
\end{choiceshoriz}
\end{question}
}

\element{proportions}{
\begin{question}{prop-echelle-3}
Sur un plan à l'échelle $1:3000$, $4$ cm représentent quelle distance réelle ?
\begin{choiceshoriz}
\correctchoice{$12000$ cm}
\wrongchoice{$3004$ cm}
\wrongchoice{$120$ cm}
\wrongchoice{$750$ cm}
\end{choiceshoriz}
\end{question}
}

\element{proportions}{
\begin{question}{prop-echelle-4}
Sur un plan à l'échelle $1:4000$, $5$ cm représentent quelle distance réelle ?
\begin{choiceshoriz}
\correctchoice{$20000$ cm}
\wrongchoice{$4005$ cm}
\wrongchoice{$200$ cm}
\wrongchoice{$800$ cm}
\end{choiceshoriz}
\end{question}
}

\element{proportions}{
\begin{question}{prop-echelle-5}
Sur un plan à l'échelle $1:5000$, $6$ cm représentent quelle distance réelle ?
\begin{choiceshoriz}
\correctchoice{$30000$ cm}
\wrongchoice{$5006$ cm}
\wrongchoice{$300$ cm}
\wrongchoice{$833{,}33$ cm}
\end{choiceshoriz}
\end{question}
}

\element{proportions}{
\begin{question}{prop-echelle-6}
Sur un plan à l'échelle $1:6000$, $7$ cm représentent quelle distance réelle ?
\begin{choiceshoriz}
\correctchoice{$42000$ cm}
\wrongchoice{$6007$ cm}
\wrongchoice{$420$ cm}
\wrongchoice{$857{,}14$ cm}
\end{choiceshoriz}
\end{question}
}

\element{proportions}{
\begin{question}{prop-echelle-7}
Sur un plan à l'échelle $1:7000$, $8$ cm représentent quelle distance réelle ?
\begin{choiceshoriz}
\correctchoice{$56000$ cm}
\wrongchoice{$7008$ cm}
\wrongchoice{$560$ cm}
\wrongchoice{$875$ cm}
\end{choiceshoriz}
\end{question}
}

\element{proportions}{
\begin{question}{prop-echelle-8}
Sur un plan à l'échelle $1:8000$, $9$ cm représentent quelle distance réelle ?
\begin{choiceshoriz}
\correctchoice{$72000$ cm}
\wrongchoice{$8009$ cm}
\wrongchoice{$720$ cm}
\wrongchoice{$888{,}89$ cm}
\end{choiceshoriz}
\end{question}
}

\element{proportions}{
\begin{question}{prop-echelle-9}
Sur un plan à l'échelle $1:9000$, $10$ cm représentent quelle distance réelle ?
\begin{choiceshoriz}
\correctchoice{$90000$ cm}
\wrongchoice{$9010$ cm}
\wrongchoice{$900$ cm}
\wrongchoice{$1$}
\end{choiceshoriz}
\end{question}
}

\element{proportions}{
\begin{question}{prop-echelle-10}
Sur un plan à l'échelle $1:10000$, $11$ cm représentent quelle distance réelle ?
\begin{choiceshoriz}
\correctchoice{$110000$ cm}
\wrongchoice{$10011$ cm}
\wrongchoice{$1100$ cm}
\wrongchoice{$909{,}09$ cm}
\end{choiceshoriz}
\end{question}
}

\element{proportions}{
\begin{question}{prop-fraction-pourcentage-1}
Convertir $\frac{1}{2}$ en pourcentage.
\begin{choiceshoriz}
\correctchoice{$50\%$}
\wrongchoice{$0{,}5\%$}
\wrongchoice{$2\%$}
\wrongchoice{$200\%$}
\end{choiceshoriz}
\end{question}
}

\element{proportions}{
\begin{question}{prop-fraction-pourcentage-2}
Convertir $\frac{1}{4}$ en pourcentage.
\begin{choiceshoriz}
\correctchoice{$25\%$}
\wrongchoice{$0{,}25\%$}
\wrongchoice{$4\%$}
\wrongchoice{$400\%$}
\end{choiceshoriz}
\end{question}
}

\element{proportions}{
\begin{question}{prop-fraction-pourcentage-3}
Convertir $\frac{1}{5}$ en pourcentage.
\begin{choiceshoriz}
\correctchoice{$20\%$}
\wrongchoice{$0{,}2\%$}
\wrongchoice{$5\%$}
\wrongchoice{$500\%$}
\end{choiceshoriz}
\end{question}
}

\element{proportions}{
\begin{question}{prop-fraction-pourcentage-4}
Convertir $\frac{2}{8}$ en pourcentage.
\begin{choiceshoriz}
\correctchoice{$25\%$}
\wrongchoice{$0{,}25\%$}
\wrongchoice{$4\%$}
\wrongchoice{$400\%$}
\end{choiceshoriz}
\end{question}
}

\element{proportions}{
\begin{question}{prop-fraction-pourcentage-5}
Convertir $\frac{2}{10}$ en pourcentage.
\begin{choiceshoriz}
\correctchoice{$20\%$}
\wrongchoice{$0{,}2\%$}
\wrongchoice{$5\%$}
\wrongchoice{$500\%$}
\end{choiceshoriz}
\end{question}
}

\element{proportions}{
\begin{question}{prop-fraction-pourcentage-6}
Convertir $\frac{5}{20}$ en pourcentage.
\begin{choiceshoriz}
\correctchoice{$25\%$}
\wrongchoice{$0{,}25\%$}
\wrongchoice{$4\%$}
\wrongchoice{$400\%$}
\end{choiceshoriz}
\end{question}
}

\element{proportions}{
\begin{question}{prop-fraction-pourcentage-7}
Convertir $\frac{6}{25}$ en pourcentage.
\begin{choiceshoriz}
\correctchoice{$24\%$}
\wrongchoice{$0{,}24\%$}
\wrongchoice{$4{,}17\%$}
\wrongchoice{$416{,}67\%$}
\end{choiceshoriz}
\end{question}
}

\element{proportions}{
\begin{question}{prop-fraction-pourcentage-8}
Convertir $\frac{10}{40}$ en pourcentage.
\begin{choiceshoriz}
\correctchoice{$25\%$}
\wrongchoice{$0{,}25\%$}
\wrongchoice{$4\%$}
\wrongchoice{$400\%$}
\end{choiceshoriz}
\end{question}
}

\element{proportions}{
\begin{question}{prop-fraction-pourcentage-9}
Convertir $\frac{12}{50}$ en pourcentage.
\begin{choiceshoriz}
\correctchoice{$24\%$}
\wrongchoice{$0{,}24\%$}
\wrongchoice{$4{,}17\%$}
\wrongchoice{$416{,}67\%$}
\end{choiceshoriz}
\end{question}
}

\element{proportions}{
\begin{question}{prop-fraction-pourcentage-10}
Convertir $\frac{25}{100}$ en pourcentage.
\begin{choiceshoriz}
\correctchoice{$25\%$}
\wrongchoice{$0{,}25\%$}
\wrongchoice{$4\%$}
\wrongchoice{$400\%$}
\end{choiceshoriz}
\end{question}
}

\element{proportions}{
\begin{question}{prop-coeff-hausse-1}
Quel coefficient multiplicateur correspond à une hausse de $5\%$ ?
\begin{choiceshoriz}
\correctchoice{$1{,}05$}
\wrongchoice{$0{,}05$}
\wrongchoice{$0{,}95$}
\wrongchoice{$105$}
\end{choiceshoriz}
\end{question}
}

\element{proportions}{
\begin{question}{prop-coeff-hausse-2}
Quel coefficient multiplicateur correspond à une hausse de $10\%$ ?
\begin{choiceshoriz}
\correctchoice{$1{,}1$}
\wrongchoice{$0{,}1$}
\wrongchoice{$0{,}9$}
\wrongchoice{$110$}
\end{choiceshoriz}
\end{question}
}

\element{proportions}{
\begin{question}{prop-coeff-hausse-3}
Quel coefficient multiplicateur correspond à une hausse de $15\%$ ?
\begin{choiceshoriz}
\correctchoice{$1{,}15$}
\wrongchoice{$0{,}15$}
\wrongchoice{$0{,}85$}
\wrongchoice{$115$}
\end{choiceshoriz}
\end{question}
}

\element{proportions}{
\begin{question}{prop-coeff-hausse-4}
Quel coefficient multiplicateur correspond à une hausse de $20\%$ ?
\begin{choiceshoriz}
\correctchoice{$1{,}2$}
\wrongchoice{$0{,}2$}
\wrongchoice{$0{,}8$}
\wrongchoice{$120$}
\end{choiceshoriz}
\end{question}
}

\element{proportions}{
\begin{question}{prop-coeff-hausse-5}
Quel coefficient multiplicateur correspond à une hausse de $25\%$ ?
\begin{choiceshoriz}
\correctchoice{$1{,}25$}
\wrongchoice{$0{,}25$}
\wrongchoice{$0{,}75$}
\wrongchoice{$125$}
\end{choiceshoriz}
\end{question}
}

\element{proportions}{
\begin{question}{prop-coeff-hausse-6}
Quel coefficient multiplicateur correspond à une hausse de $30\%$ ?
\begin{choiceshoriz}
\correctchoice{$1{,}3$}
\wrongchoice{$0{,}3$}
\wrongchoice{$0{,}7$}
\wrongchoice{$130$}
\end{choiceshoriz}
\end{question}
}

\element{proportions}{
\begin{question}{prop-coeff-hausse-7}
Quel coefficient multiplicateur correspond à une hausse de $35\%$ ?
\begin{choiceshoriz}
\correctchoice{$1{,}35$}
\wrongchoice{$0{,}35$}
\wrongchoice{$0{,}65$}
\wrongchoice{$135$}
\end{choiceshoriz}
\end{question}
}

\element{proportions}{
\begin{question}{prop-coeff-hausse-8}
Quel coefficient multiplicateur correspond à une hausse de $40\%$ ?
\begin{choiceshoriz}
\correctchoice{$1{,}4$}
\wrongchoice{$0{,}4$}
\wrongchoice{$0{,}6$}
\wrongchoice{$140$}
\end{choiceshoriz}
\end{question}
}

\element{proportions}{
\begin{question}{prop-coeff-hausse-9}
Quel coefficient multiplicateur correspond à une hausse de $45\%$ ?
\begin{choiceshoriz}
\correctchoice{$1{,}45$}
\wrongchoice{$0{,}45$}
\wrongchoice{$0{,}55$}
\wrongchoice{$145$}
\end{choiceshoriz}
\end{question}
}

\element{proportions}{
\begin{question}{prop-coeff-hausse-10}
Quel coefficient multiplicateur correspond à une hausse de $50\%$ ?
\begin{choiceshoriz}
\correctchoice{$1{,}5$}
\wrongchoice{$0{,}5$}
\wrongchoice{$150$}
\wrongchoice{$1$}
\end{choiceshoriz}
\end{question}
}

\element{proportions}{
\begin{question}{prop-comparaison-1}
Quelle part représente $20$ sur $50$ ?
\begin{choiceshoriz}
\correctchoice{$40\%$}
\wrongchoice{$2{,}5\%$}
\wrongchoice{$70\%$}
\wrongchoice{$250\%$}
\end{choiceshoriz}
\end{question}
}

\element{proportions}{
\begin{question}{prop-comparaison-2}
Quelle part représente $25$ sur $60$ ?
\begin{choiceshoriz}
\correctchoice{$41{,}67\%$}
\wrongchoice{$2{,}4\%$}
\wrongchoice{$85\%$}
\wrongchoice{$240\%$}
\end{choiceshoriz}
\end{question}
}

\element{proportions}{
\begin{question}{prop-comparaison-3}
Quelle part représente $30$ sur $70$ ?
\begin{choiceshoriz}
\correctchoice{$42{,}86\%$}
\wrongchoice{$2{,}33\%$}
\wrongchoice{$100\%$}
\wrongchoice{$233{,}33\%$}
\end{choiceshoriz}
\end{question}
}

\element{proportions}{
\begin{question}{prop-comparaison-4}
Quelle part représente $35$ sur $80$ ?
\begin{choiceshoriz}
\correctchoice{$43{,}75\%$}
\wrongchoice{$2{,}29\%$}
\wrongchoice{$115\%$}
\wrongchoice{$228{,}57\%$}
\end{choiceshoriz}
\end{question}
}

\element{proportions}{
\begin{question}{prop-comparaison-5}
Quelle part représente $40$ sur $90$ ?
\begin{choiceshoriz}
\correctchoice{$44{,}44\%$}
\wrongchoice{$2{,}25\%$}
\wrongchoice{$130\%$}
\wrongchoice{$225\%$}
\end{choiceshoriz}
\end{question}
}

\element{proportions}{
\begin{question}{prop-comparaison-6}
Quelle part représente $45$ sur $100$ ?
\begin{choiceshoriz}
\correctchoice{$45\%$}
\wrongchoice{$2{,}22\%$}
\wrongchoice{$145\%$}
\wrongchoice{$222{,}22\%$}
\end{choiceshoriz}
\end{question}
}

\element{proportions}{
\begin{question}{prop-comparaison-7}
Quelle part représente $50$ sur $110$ ?
\begin{choiceshoriz}
\correctchoice{$45{,}45\%$}
\wrongchoice{$2{,}2\%$}
\wrongchoice{$160\%$}
\wrongchoice{$220\%$}
\end{choiceshoriz}
\end{question}
}

\element{proportions}{
\begin{question}{prop-comparaison-8}
Quelle part représente $55$ sur $120$ ?
\begin{choiceshoriz}
\correctchoice{$45{,}83\%$}
\wrongchoice{$2{,}18\%$}
\wrongchoice{$175\%$}
\wrongchoice{$218{,}18\%$}
\end{choiceshoriz}
\end{question}
}

\element{proportions}{
\begin{question}{prop-comparaison-9}
Quelle part représente $60$ sur $130$ ?
\begin{choiceshoriz}
\correctchoice{$46{,}15\%$}
\wrongchoice{$2{,}17\%$}
\wrongchoice{$190\%$}
\wrongchoice{$216{,}67\%$}
\end{choiceshoriz}
\end{question}
}

\element{proportions}{
\begin{question}{prop-comparaison-10}
Quelle part représente $65$ sur $140$ ?
\begin{choiceshoriz}
\correctchoice{$46{,}43\%$}
\wrongchoice{$2{,}15\%$}
\wrongchoice{$205\%$}
\wrongchoice{$215{,}38\%$}
\end{choiceshoriz}
\end{question}
}
